\documentclass[12pt,a4paper]{article}

\usepackage[a4paper,top=2.3cm,bottom=2.3cm,left=2.6cm,right=2.4cm,headheight=17pt]{geometry}
\usepackage{fontspec}
\setmainfont{DejaVu Serif}
\setsansfont{DejaVu Sans}
\setmonofont{DejaVu Sans Mono}

\usepackage{amsmath,amssymb}
\usepackage{xcolor}
\definecolor{FPTBlue}{HTML}{1F4E79}
\usepackage{graphicx}
\usepackage{float}
\usepackage{caption}
\usepackage{longtable,booktabs,array,calc}
\usepackage{ragged2e}
\usepackage{enumitem}
\usepackage{setspace}
\usepackage{microtype}
\usepackage{titlesec}
\usepackage{fancyhdr}
\usepackage{xurl}
\usepackage[hidelinks,unicode]{hyperref}
\usepackage{bookmark}
\usepackage{etoolbox}

\hypersetup{
  pdftitle={Báo cáo cuối kỳ AIL303m - Nhận diện biển báo giao thông Việt Nam},
  pdfsubject={HOG, HSV Histogram, PCA và Linear SVM},
  pdfcreator={XeLaTeX},
  colorlinks=false
}

\renewcommand{\contentsname}{MỤC LỤC}
\setcounter{secnumdepth}{-1}
\setcounter{tocdepth}{2}
\setlength{\parindent}{1.2cm}
\setlength{\parskip}{0.25em}
\setlength{\emergencystretch}{4em}
\onehalfspacing
\justifying

\titleformat{\section}
  {\Large\bfseries\color{FPTBlue}}{}{0pt}{}
\titleformat{\subsection}
  {\large\bfseries\color{FPTBlue}}{}{0pt}{}
\titlespacing*{\section}{0pt}{1.5em}{0.7em}
\titlespacing*{\subsection}{0pt}{1.1em}{0.45em}

\setlist[itemize]{leftmargin=1.0cm,itemsep=0.2em,topsep=0.3em}
\setlist[enumerate]{leftmargin=1.0cm,itemsep=0.2em,topsep=0.3em}
\AtBeginEnvironment{longtable}{\small\setlength{\tabcolsep}{4pt}}
\captionsetup{font=small,labelfont=bf,justification=centering}

\pagestyle{fancy}
\fancyhf{}
\fancyhead[C]{\small\itshape BÁO CÁO CUỐI KỲ - NHẬN DIỆN BIỂN BÁO GIAO THÔNG VIỆT NAM}
\fancyfoot[C]{\thepage}
\renewcommand{\headrulewidth}{0.4pt}
\renewcommand{\footrulewidth}{0pt}

% Pandoc compatibility helpers.
\makeatletter
\patchcmd\longtable{\par}{\if@noskipsec\mbox{}\fi\par}{}{}
\makeatother
\providecommand{\tightlist}{\setlength{\itemsep}{0pt}\setlength{\parskip}{0pt}}
\makeatletter
\def\maxwidth{\ifdim\Gin@nat@width>\linewidth\linewidth\else\Gin@nat@width\fi}
\def\maxheight{\ifdim\Gin@nat@height>0.82\textheight 0.82\textheight\else\Gin@nat@height\fi}
\makeatother
\setkeys{Gin}{width=\maxwidth,height=\maxheight,keepaspectratio}
\urlstyle{tt}
\sloppy

\begin{document}
\begin{titlepage}
  \centering
  {\large\bfseries TRƯỜNG ĐẠI HỌC FPT\par}
  \vspace{0.25cm}
  {\large\bfseries FPT UNIVERSITY ĐÀ NẴNG\par}
  \vfill
  {\Large\bfseries BÁO CÁO CUỐI KỲ\par}
  \vspace{0.8cm}
  {\LARGE\bfseries NHẬN DIỆN BIỂN BÁO GIAO THÔNG VIỆT NAM\par}
  \vspace{0.35cm}
  {\Large\bfseries BẰNG CÁC PHƯƠNG PHÁP TRÍCH XUẤT ĐẶC TRƯNG VÀ MÔ HÌNH HỌC MÁY\par}
  \vfill
  \begin{tabular}{@{}p{0.34\textwidth}p{0.54\textwidth}@{}}
    \textbf{Môn học:} & AIL303m \\
    \textbf{Học kỳ:} & Summer 2026 \\
    \textbf{Giảng viên hướng dẫn:} & 	Võ Quốc Trình \\
    \textbf{Nhóm thực hiện:} & 6\\
    \textbf{Danh sách thành viên:} & Lê Khôi Nguyên (SE200064), 
    Bùi	Anh	Kiệt (DE200465), 
    Trương Phước Nguyên (DE190562), 
    Phan Hữu Nhật Trường (DE200114) \\
  \end{tabular}
  \vfill
  {\large Đà Nẵng, tháng 7 năm 2026\par}
\end{titlepage}

\section*{TÓM TẮT}\addcontentsline{toc}{section}{TÓM TẮT}

Báo cáo trình bày quá trình xây dựng, kiểm chứng và tối ưu một hệ thống
phân loại biển báo giao thông Việt Nam trên bộ dữ liệu NVTS gồm 47 lớp.
Bài toán có ba thách thức nổi bật: nhiều lớp có hình dạng và màu sắc gần
giống nhau, ảnh thực tế chịu ảnh hưởng của mờ, chiếu sáng và biến dạng
góc nhìn, và phân bố lớp mất cân bằng nghiêm trọng khi 10 lớp có dưới 30
ảnh train. Mục tiêu của dự án là đồng thời đạt hiệu năng phân loại cao,
cải thiện khả năng nhận diện lớp hiếm và duy trì độ trễ thấp trên nền
tảng học máy cổ điển, không phụ thuộc vào mô hình học sâu cồng kềnh.
Pipeline cuối cùng chuẩn hóa ảnh về 64×64 bằng pad\_square; trích xuất
HOG trên ảnh xám để mô tả hình dạng và cạnh; xây dựng Color Histogram
trong không gian HSV để mô tả màu; dung hợp hai đặc trưng; chuẩn hóa
bằng StandardScaler; giảm chiều bằng PCA giữ 95\% phương sai; tăng cường
dữ liệu trực tiếp trên ảnh cho các lớp thiểu số; và phân loại bằng SVC
kernel tuyến tính. Sau các ablation về kích thước ảnh, không gian màu,
PCA và augmentation, dự án benchmark sáu họ mô hình và khảo sát các
kernel SVM. Cấu hình cuối là SVC(kernel="linear", C=0,1,
class\_weight=None). Trên tập test 851 ảnh, mô hình đạt Accuracy
96,12\%, Macro Precision 97,51\%, Macro Recall 95,62\%, Macro F1 96,29\%
và Weighted F1 96,09\%. Thời gian 1,52 ms/ảnh là độ trễ của bước phân
loại trên vector đặc trưng đã được xử lý, không phải latency end-to-end.
So với RBF ở Phase 1, Linear SVM tăng Macro F1 0,65 điểm phần trăm và
Macro Recall 1,23 điểm phần trăm, đồng thời giảm 817 support vector. Kết
quả cho thấy một pipeline đặc trưng được thiết kế phù hợp có thể tạo ra
không gian phân loại chất lượng cao, trong đó mô hình tuyến tính vừa
cạnh tranh về hiệu năng vừa thuận lợi cho triển khai tài nguyên hạn chế.

\textbf{Từ khóa:} \emph{nhận diện biển báo giao thông, HOG, HSV
Histogram, PCA, mất cân bằng lớp, Linear SVM, học máy truyền thống.}

\begin{longtable}[]{@{}
  >{\raggedright\arraybackslash}p{(\columnwidth - 0\tabcolsep) * \real{1.0000}}@{}}
\toprule\noalign{}
\begin{minipage}[b]{\linewidth}\raggedright
\textbf{Thông điệp chính}

Kết quả tốt nhất không đến từ việc tăng độ phức tạp của mô hình, mà đến
từ sự phối hợp giữa biểu diễn hình dạng, màu sắc, giảm chiều và quy
trình lựa chọn mô hình dựa trên Macro F1.
\end{minipage} \\
\midrule\noalign{}
\endhead
\bottomrule\noalign{}
\endlastfoot
\end{longtable}

\clearpage
\tableofcontents
\clearpage
\section*{DANH MỤC TỪ VIẾT TẮT}\addcontentsline{toc}{section}{DANH MỤC TỪ VIẾT TẮT}

\begin{longtable}[]{@{}
  >{\raggedright\arraybackslash}p{(\columnwidth - 2\tabcolsep) * \real{0.5000}}
  >{\raggedright\arraybackslash}p{(\columnwidth - 2\tabcolsep) * \real{0.5000}}@{}}
\toprule\noalign{}
\begin{minipage}[b]{\linewidth}\raggedright
\textbf{Viết tắt}
\end{minipage} & \begin{minipage}[b]{\linewidth}\raggedright
\textbf{Nội dung}
\end{minipage} \\
\midrule\noalign{}
\endhead
\bottomrule\noalign{}
\endlastfoot
ADAS & Advanced Driver Assistance Systems - Hệ thống hỗ trợ người lái
nâng cao \\
HOG & Histogram of Oriented Gradients \\
HSV & Hue - Saturation - Value \\
PCA & Principal Component Analysis \\
SVM/SVC & Support Vector Machine / Support Vector Classifier \\
RBF & Radial Basis Function \\
KNN & K-Nearest Neighbors \\
MLP & Multi-layer Perceptron \\
NVTS & Vietnamese Traffic Sign Dataset \\
\end{longtable}

\section{1. GIỚI THIỆU}\label{giux1edbi-thiux1ec7u}

\subsection{1.1. Bối cảnh nghiên
cứu}\label{bux1ed1i-cux1ea3nh-nghiuxean-cux1ee9u}

Nhận diện biển báo giao thông là một thành phần quan trọng trong các hệ
thống hỗ trợ người lái, giám sát giao thông và nền tảng giao thông thông
minh. Một hệ thống đáng tin cậy phải phân biệt được nhiều loại biển có
hình dạng, màu sắc và biểu tượng gần giống nhau, đồng thời duy trì tốc
độ xử lý đủ nhanh trong điều kiện ảnh có thể bị mờ, thiếu sáng, lệch
phối cảnh, che khuất hoặc suy giảm độ phân giải. Các nghiên cứu trước
cho thấy đặc trưng hình dạng như HOG có khả năng biểu diễn tốt cấu trúc
cạnh của biển báo, trong khi thông tin màu giúp phân biệt các nhóm biển
cấm, cảnh báo và chỉ dẫn {[}2{]}, {[}3{]}.

Học sâu hiện là hướng chủ đạo trong nhiều bài toán thị giác máy tính,
nhưng thường cần lượng dữ liệu gán nhãn lớn, tài nguyên huấn luyện cao
và quá trình giải thích mô hình phức tạp hơn. Trong phạm vi môn AIL303m
và mục tiêu xây dựng một hệ thống nhẹ, dự án lựa chọn classical computer
vision kết hợp machine learning. Cách tiếp cận này cho phép đánh giá
riêng tác động của tiền xử lý, đặc trưng, giảm chiều và bộ phân loại;
đồng thời phù hợp với các môi trường không có GPU.

\subsection{1.2. Phát biểu bài toán và phạm
vi}\label{phuxe1t-biux1ec3u-buxe0i-touxe1n-vuxe0-phux1ea1m-vi}

Đầu vào của hệ thống là một ảnh đã được crop, trong đó mỗi ảnh chứa một
biển báo giao thông. Nhiệm vụ là gán ảnh vào một trong 47 lớp của bộ dữ
liệu NVTS. Phạm vi hiện tại là phân loại biển báo, chưa bao gồm bước
phát hiện vị trí biển trong ảnh đường phố hoàn chỉnh. Vì vậy, kết quả
của báo cáo phản ánh năng lực nhận diện trên ảnh đã crop và không được
diễn giải như hiệu năng end-to-end của một hệ thống phát hiện - nhận
diện hoàn chỉnh.

Dữ liệu gồm 6.605 ảnh train, 824 ảnh validation và 851 ảnh test. Thách
thức quan trọng là mất cân bằng lớp: 10 lớp có dưới 30 ảnh train, trong
đó R.301e chỉ có 16 ảnh và W.205c có 17 ảnh. Nếu chỉ tối ưu Accuracy, mô
hình có thể ưu tiên các lớp phổ biến và vẫn đạt tỷ lệ đúng cao dù bỏ sót
lớp hiếm. Vì vậy, Macro F1 được chọn làm metric chính.

\subsection{1.3. Mục tiêu nghiên
cứu}\label{mux1ee5c-tiuxeau-nghiuxean-cux1ee9u}

\begin{itemize}
\item
  Tìm cấu hình tiền xử lý ảnh giữ được hình học đặc trưng của biển báo
  và có chi phí tính toán hợp lý.
\item
  Xây dựng biểu diễn kết hợp thông tin hình dạng và màu sắc bằng HOG
  Gray và HSV Color Histogram.
\item
  Đánh giá tác động của StandardScaler và PCA đối với hiệu năng, số
  chiều và độ trễ.
\item
  Cải thiện khả năng nhận diện lớp hiếm bằng tăng cường dữ liệu trực
  tiếp trên ảnh train.
\item
  Benchmark nhiều họ mô hình học máy và tinh chỉnh SVM để tìm cấu hình
  cuối cùng tốt nhất.
\item
  Đánh giá mô hình bằng Accuracy, Macro Precision, Macro Recall, Macro
  F1, Weighted F1 và thời gian suy luận.
\end{itemize}

\subsection{1.4. Câu hỏi nghiên
cứu}\label{cuxe2u-hux1ecfi-nghiuxean-cux1ee9u}

\begin{enumerate}
\def\labelenumi{\arabic{enumi}.}
\item
  Kết hợp HOG và HSV Histogram có tạo ra biểu diễn hiệu quả hơn khi xét
  trong pipeline hoàn chỉnh hay không?
\item
  PCA có thể giảm mạnh số chiều mà vẫn duy trì hoặc cải thiện Macro F1
  hay không?
\item
  Tăng cường dữ liệu cho lớp hiếm có cải thiện Macro Recall và Macro F1
  hay không?
\item
  Sau khi đặc trưng được chuẩn hóa và giảm chiều, mô hình tuyến tính hay
  phi tuyến phù hợp hơn?
\item
  Độ trễ được báo cáo có ý nghĩa như thế nào đối với khả năng triển khai
  thời gian thực?
\end{enumerate}

\subsection{1.5. Đóng góp chính}\label{ux111uxf3ng-guxf3p-chuxednh}

\begin{enumerate}
\def\labelenumi{\arabic{enumi}.}
\setcounter{enumi}{5}
\item
  Một pipeline hoàn chỉnh gồm 64×64 pad\_square, HOG Gray, HSV
  Histogram, StandardScaler, PCA 95\%, minority augmentation và Linear
  SVM.
\item
  Một quy trình thực nghiệm hai giai đoạn: Phase 1 khóa mô hình để đánh
  giá feature engineering; Phase 2 benchmark và tuning mô hình.
\item
  Một giải pháp xử lý lớp hiếm bằng augmentation ảnh 2D, bổ sung 430 ảnh
  train mà không nội suy trực tiếp trong không gian HOG.
\item
  Một phân tích đầy đủ về trade-off giữa Accuracy, Macro Precision,
  Macro Recall, Macro F1, độ trễ và số support vector.
\item
  Một phần audit khoa học chỉ rõ giới hạn của protocol test, latency và
  sự chưa thống nhất của số chiều đặc trưng trong các artifact.
\end{enumerate}

\section{2. CƠ SỞ LÝ THUYẾT VÀ NGHIÊN CỨU LIÊN
QUAN}\label{cux1a1-sux1edf-luxfd-thuyux1ebft-vuxe0-nghiuxean-cux1ee9u-liuxean-quan}

\subsection{2.1. Pipeline nhận diện biển báo bằng đặc trưng thủ
công}\label{pipeline-nhux1eadn-diux1ec7n-biux1ec3n-buxe1o-bux1eb1ng-ux111ux1eb7c-trux1b0ng-thux1ee7-cuxf4ng}

Pipeline truyền thống cho nhận diện biển báo thường gồm tiền xử lý,
trích xuất đặc trưng và phân loại. Chất lượng biểu diễn ảnh quyết định
đáng kể đến khả năng phân tách của mô hình. Với biển báo, hai nguồn
thông tin quan trọng là hình dạng - cạnh và màu sắc. HOG mô tả cấu trúc
gradient; histogram màu mô tả phân bố sắc độ; bộ phân loại sau đó hoạt
động trên vector đặc trưng thay vì pixel thô.

\subsection{2.2. Histogram of Oriented
Gradients}\label{histogram-of-oriented-gradients}

HOG do Dalal và Triggs giới thiệu để mô tả phân bố hướng gradient trong
các vùng cục bộ {[}2{]}. Ảnh được chia thành các cell; mỗi cell tích lũy
histogram của hướng gradient, có trọng số theo độ lớn gradient; các cell
lân cận được chuẩn hóa theo block. HOG bền vững hơn pixel thô trước một
số thay đổi chiếu sáng và đặc biệt phù hợp với biển báo vì chúng có
đường viền mạnh, hình học rõ và biểu tượng tương phản cao.

Trong dự án, HOG được trích xuất trên ảnh grayscale với orientations=9,
pixels\_per\_cell=(8,8), cells\_per\_block=(2,2) và block\_norm=L2-Hys.
Các thông số này được dùng làm mô tả cấu hình; tuy nhiên số chiều tổng
1.812 trong artifact M04 chưa khớp hoàn toàn với phép tính lý thuyết từ
các tham số được ghi lại. Báo cáo giữ nguyên giá trị artifact và đưa vấn
đề này vào phần hạn chế thay vì tự suy diễn.

\subsection{2.3. Color Histogram trong không gian
HSV}\label{color-histogram-trong-khuxf4ng-gian-hsv}

Không gian HSV tách sắc độ Hue khỏi độ bão hòa Saturation và độ sáng
Value. Điều này thuận lợi khi mô tả màu biển báo vì các nhóm biển thường
tuân theo quy ước màu tương đối ổn định: đỏ cho nhiều biển cấm, vàng cho
cảnh báo, xanh cho chỉ dẫn. Histogram 3D với 8 bin mỗi kênh tạo 512 bin
và được chuẩn hóa theo tổng khối lượng. Histogram không giữ vị trí không
gian chi tiết, nhưng cung cấp tín hiệu màu tổng thể bổ sung cho HOG.

\subsection{2.4. PCA và giảm
chiều}\label{pca-vuxe0-giux1ea3m-chiux1ec1u}

Khi nhiều đặc trưng được nối với nhau, số chiều tăng nhanh và có thể
chứa nhiều thành phần tương quan hoặc dư thừa. PCA biến đổi dữ liệu sang
hệ trục trực giao được sắp theo phương sai giảm dần. Trong dự án,
StandardScaler được fit trên train trước PCA; PCA giữ 95\% phương sai và
giảm vector cuối xuống 571 chiều. PCA không trực tiếp tối ưu nhãn lớp và
không phải lúc nào cũng là bộ lọc nhiễu; hiệu quả của nó phải được xác
nhận bằng thực nghiệm.

\subsection{2.5. Support Vector Machine}\label{support-vector-machine}

SVM tìm siêu phẳng có lề phân cách lớn giữa các lớp. Với soft margin, mô
hình cho phép một số điểm vi phạm lề để cải thiện khả năng khái quát.
Bài toán nhị phân có thể viết dưới dạng:

\begin{equation}                                         
\min_{w, b, \xi} \left( \frac{1}{2} \|w\|^2 + C      
\sum_{i=1}^{n} \xi_i \right)                               
\end{equation}                                           
với các ràng buộc:                                       
\begin{equation}                                         
y_i (w^{\mathsf{T}}x_i + b) \ge 1 - \xi_i, \quad     
\text{và} \quad \xi_i \ge 0, \quad \forall i = 1, 2, \dots,
n                                                          
\end{equation}   

Tham số C điều khiển đánh đổi giữa lề rộng và mức phạt vi phạm. Kernel
tuyến tính phù hợp khi dữ liệu có thể phân tách tốt trong không gian
hiện tại; RBF cho ranh giới linh hoạt hơn nhưng đòi hỏi tuning C và
gamma. Với bài toán đa lớp, SVC xây dựng nhiều bài toán nhị phân nội bộ.
Trong dự án, cấu hình Linear SVM với C=0,1 cho validation Macro F1 cao
nhất.

\subsection{2.6. Đánh giá trên dữ liệu mất cân
bằng}\label{ux111uxe1nh-giuxe1-truxean-dux1eef-liux1ec7u-mux1ea5t-cuxe2n-bux1eb1ng}

Accuracy đo tỷ lệ dự đoán đúng trên toàn bộ mẫu, nhưng có thể che khuất
hiệu năng kém của lớp hiếm. Macro Precision, Macro Recall và Macro F1
tính metric riêng cho từng lớp rồi lấy trung bình không trọng số; do đó
mỗi lớp đóng góp ngang nhau. Weighted F1 được báo cáo bổ sung để phản
ánh hiệu năng theo phân bố mẫu thực tế. Dự án chọn Macro F1 làm tiêu chí
chính để cân bằng giữa precision và recall trên 47 lớp.

\subsection{2.7. Nghiên cứu liên
quan}\label{nghiuxean-cux1ee9u-liuxean-quan}

Ma và Huang đề xuất hệ thống nhận diện biển báo phân cấp coarse-to-fine,
kết hợp HOG ở tầng thô và Dense-SIFT, LBP, Gabor cùng committee SVM -
Random Forest ở tầng chi tiết. Phương pháp đạt 98,76\% Accuracy trên
GTSRB và khoảng 50 ms/ảnh trên phần cứng được báo cáo {[}3{]}. Công
trình này củng cố hai nhận định: đặc trưng bổ sung có thể tăng hiệu năng
và chiến lược phân cấp phù hợp với các lớp có ngoại hình gần giống nhau.
Tuy nhiên, không thể so sánh trực tiếp con số 98,76\% với kết quả NVTS
của dự án vì bộ dữ liệu, số lớp, kích thước train/test, protocol và phần
cứng khác nhau.

GTSRB là benchmark lớn với 43 lớp và hơn 50 nghìn ảnh, thường được dùng
để so sánh các hệ thống nhận diện biển báo {[}4{]}. Dự án hiện tại tập
trung vào NVTS 47 lớp và mục tiêu lightweight classical ML. Vì vậy,
related work được dùng để định hướng thiết kế và thảo luận, không dùng
để khẳng định superiority tuyệt đối.

\section{3. DỮ LIỆU VÀ PHÂN TÍCH MẤT CÂN
BẰNG}\label{dux1eef-liux1ec7u-vuxe0-phuxe2n-tuxedch-mux1ea5t-cuxe2n-bux1eb1ng}

\subsection{3.1. Bộ dữ liệu NVTS}\label{bux1ed9-dux1eef-liux1ec7u-nvts}

Dự án sử dụng bộ dữ liệu biển báo giao thông Việt Nam NVTS gồm 47 lớp.
Ảnh đã được crop theo biển báo và được tổ chức thành train, validation
và test. Train được dùng để fit scaler, PCA và mô hình; validation dùng
để lựa chọn cấu hình và siêu tham số; test dùng cho các đánh giá được
artifact ghi nhận. Sau tăng cường, train tăng từ 6.605 lên 7.035 ảnh.

\begin{longtable}[]{@{}
  >{\raggedright\arraybackslash}p{(\columnwidth - 6\tabcolsep) * \real{0.2500}}
  >{\raggedright\arraybackslash}p{(\columnwidth - 6\tabcolsep) * \real{0.2500}}
  >{\raggedright\arraybackslash}p{(\columnwidth - 6\tabcolsep) * \real{0.2500}}
  >{\raggedright\arraybackslash}p{(\columnwidth - 6\tabcolsep) * \real{0.2500}}@{}}
\toprule\noalign{}
\begin{minipage}[b]{\linewidth}\raggedright
\textbf{Tập dữ liệu}
\end{minipage} & \begin{minipage}[b]{\linewidth}\raggedright
\textbf{Số lớp}
\end{minipage} & \begin{minipage}[b]{\linewidth}\raggedright
\textbf{Số ảnh}
\end{minipage} & \begin{minipage}[b]{\linewidth}\raggedright
\textbf{Vai trò}
\end{minipage} \\
\midrule\noalign{}
\endhead
\bottomrule\noalign{}
\endlastfoot
Train gốc & 47 & 6.605 & Huấn luyện bộ biến đổi và mô hình \\
Validation & 47 & 824 & So sánh cấu hình và chọn siêu tham số \\
Test & 47 & 851 & Đánh giá theo các artifact dự án \\
Train sau tăng cường & 47 & 7.035 & Train cuối sau bổ sung 430 ảnh lớp
hiếm \\
\end{longtable}

\begin{center}\small\textit{Bảng 1. Quy mô và vai trò của các tập dữ liệu.}\end{center}

\begin{figure}[H]
\centering
\includegraphics[width=0.92\textwidth]{./media/image1.png}
\caption*{\textit{Hình 1. Quy mô các tập dữ liệu sử dụng trong dự án.}}
\end{figure}

\begin{figure}[H]
\centering
\includegraphics[width=0.98\textwidth]{./media/image7.png}
\caption*{\textit{Hình 2. Một số mẫu ảnh biển báo giao thông Việt Nam sau khi crop từ bộ dữ liệu NVTS.}}
\end{figure}

Hình 2 cho thấy tập dữ liệu bao gồm nhiều nhóm biển khác nhau như chỉ
hướng, cấm, hiệu lệnh và cảnh báo. Các mẫu sau crop có sự khác biệt về
kích thước vùng biển trong khung hình, độ sắc nét, ánh sáng, góc nhìn và
mức độ nhiễu nền. Điều này giải thích vì sao pipeline cần kết hợp đồng
thời thông tin hình dạng và màu sắc, thay vì chỉ dựa vào một loại đặc
trưng đơn lẻ.

\subsection{3.2. Các lớp thiểu số nghiêm
trọng}\label{cuxe1c-lux1edbp-thiux1ec3u-sux1ed1-nghiuxeam-trux1ecdng}

Phân tích M01 xác định 10 lớp có dưới 30 ảnh train. Đây là các lớp có
rủi ro cao do biến thiên trong validation/test có thể không được đại
diện đầy đủ trong train. Danh sách này được sử dụng nhất quán trong bước
augmentation ở M04.

\begin{longtable}[]{@{}
  >{\raggedright\arraybackslash}p{(\columnwidth - 2\tabcolsep) * \real{0.5000}}
  >{\raggedright\arraybackslash}p{(\columnwidth - 2\tabcolsep) * \real{0.5000}}@{}}
\toprule\noalign{}
\begin{minipage}[b]{\linewidth}\raggedright
\textbf{Nhãn lớp}
\end{minipage} & \begin{minipage}[b]{\linewidth}\raggedright
\textbf{Số ảnh train gốc}
\end{minipage} \\
\midrule\noalign{}
\endhead
\bottomrule\noalign{}
\endlastfoot
R.301e & 16 \\
W.205c & 17 \\
I.409 & 20 \\
S.505a Xe máy & 20 \\
S.505a Xe tải & 21 \\
W.239b\_ & 21 \\
P.124d & 23 \\
W.203c & 23 \\
W.205a & 26 \\
W.225 & 28 \\
\end{longtable}

\begin{center}\small\textit{Bảng 2. Mười lớp có dưới 30 ảnh train.}\end{center}

\begin{figure}[H]
\centering
\includegraphics[width=0.92\textwidth]{./media/image2.png}
\caption*{\textit{Hình 2. Phân bố số ảnh của 10 lớp thiểu số.}}
\end{figure}

Augmentation cải thiện rõ một số lớp nhưng không thể thay thế dữ liệu
thật. Đặc biệt, W.205c chỉ có 17 ảnh train và vẫn là lớp khó trong một
số đánh giá. Báo cáo không khẳng định toàn bộ các lớp hiếm đều đạt
Recall hoặc F1 cao; kết quả tổng hợp phải được đọc cùng phân tích
per-class.

\subsection{3.3. Đặc điểm và giới hạn dữ
liệu}\label{ux111ux1eb7c-ux111iux1ec3m-vuxe0-giux1edbi-hux1ea1n-dux1eef-liux1ec7u}

\begin{itemize}
\item
  Ảnh đầu vào đã crop, do đó chưa phản ánh sai số của bước phát hiện
  biển báo trong cảnh đường phố.
\item
  Một số lớp chỉ khác nhau ở chi tiết rất nhỏ như chữ số hoặc hướng mũi
  tên; độ phân giải thấp làm giảm khả năng phân biệt.
\item
  Phân bố lớp mất cân bằng khiến metric per-class của lớp hiếm có phương
  sai cao.
\item
  Hiện chỉ có một cách chia train/validation/test được báo cáo; chưa có
  repeated stratified cross-validation cho toàn pipeline.
\end{itemize}

\section{4. PHƯƠNG PHÁP ĐỀ
XUẤT}\label{phux1b0ux1a1ng-phuxe1p-ux111ux1ec1-xuux1ea5t}

\begin{figure}[H]
\centering
\includegraphics[width=0.92\textwidth]{./media/image3.png}
\caption*{\textit{Hình 3. Pipeline cuối cùng của hệ thống nhận diện.}}
\end{figure}

\subsection{4.1. Chiến lược thực nghiệm hai giai
đoạn}\label{chiux1ebfn-lux1b0ux1ee3c-thux1ef1c-nghiux1ec7m-hai-giai-ux111oux1ea1n}

Phase 1 khóa bộ phân loại ở cấu hình SVC RBF, C=10, gamma=scale và
class\_weight=balanced. Việc cố định mô hình giúp cô lập ảnh hưởng của
kích thước ảnh, chế độ resize, không gian màu, dung hợp đặc trưng, PCA
và augmentation. Sau khi chốt không gian đặc trưng, Phase 2 benchmark
nhiều thuật toán và tinh chỉnh kernel, C, gamma và class\_weight. Cách
tổ chức này hạn chế việc gán nhầm cải thiện do model tuning cho feature
engineering hoặc ngược lại.

\subsection{4.2. Tiền xử lý hình
học}\label{tiux1ec1n-xux1eed-luxfd-huxecnh-hux1ecdc}

M02 khảo sát 32×32, 48×48 và 64×64 với hai chế độ stretch và
pad\_square. Trong ablation HOG Gray độc lập, 48×48 stretch đạt
Validation Macro F1 cao nhất 95,95\%. Tuy nhiên, pipeline tích hợp
M04/M06 sử dụng 64×64 pad\_square cùng HOG, histogram màu, PCA và
augmentation. Pad\_square đệm ảnh về hình vuông trước khi resize, nhờ đó
hạn chế làm méo hình tròn, tam giác và chữ nhật. Báo cáo ghi nhận rõ sự
khác biệt giữa ``cấu hình tốt nhất trong ablation riêng'' và ``cấu hình
cuối của pipeline''.

\subsection{4.3. Đặc trưng HOG
Gray}\label{ux111ux1eb7c-trux1b0ng-hog-gray}

Ảnh được chuyển sang grayscale trước khi tính HOG. Gradient ngang và dọc
được ước lượng tại từng pixel; độ lớn và góc gradient được tổng hợp
thành histogram hướng trong các cell, sau đó chuẩn hóa theo block. Cấu
hình được ghi nhận gồm 9 hướng, pixels\_per\_cell=(8,8),
cells\_per\_block=(2,2) và block\_norm=L2-Hys. HOG tập trung vào đường
viền, góc và cấu trúc biểu tượng nên phù hợp để phân biệt hình học của
biển báo.

\subsection{4.4. Color Histogram HSV}\label{color-histogram-hsv}

Song song với HOG, ảnh màu được chuyển sang HSV và xây dựng histogram ba
chiều với 8 bin cho mỗi kênh, tạo 512 bin. Histogram được chuẩn hóa để
giảm ảnh hưởng của cường độ tổng. So với tính HOG trên nhiều kênh màu,
histogram HSV cung cấp tín hiệu màu gọn hơn và tránh mở rộng quá lớn số
chiều gradient.

\subsection{4.5. Dung hợp đặc
trưng}\label{dung-hux1ee3p-ux111ux1eb7c-trux1b0ng}

Vector HOG và HSV Histogram được nối theo chiều cột. Artifact nghiệm thu
M04 ghi nhận vector sau dung hợp có 1.812 chiều. Do giá trị này chưa
khớp hoàn toàn với phép tính lý thuyết từ bộ tham số HOG được mô tả, báo
cáo sử dụng con số 1.812 như giá trị thực nghiệm và không công bố
decomposition chi tiết khi chưa audit trực tiếp mã nguồn final. Đây là
lựa chọn thận trọng nhằm tránh tạo ra thông tin có vẻ chính xác nhưng
không được chứng minh.

\subsection{4.6. StandardScaler và PCA}\label{standardscaler-vuxe0-pca}

StandardScaler chỉ được fit trên train để đưa mỗi chiều về trung bình 0
và độ lệch chuẩn 1. Cùng scaler đó được dùng để transform validation và
test. PCA cũng chỉ fit trên train, giữ 95\% phương sai. Sau biến đổi, số
chiều pipeline cuối là 571. Quy tắc fit trên train giúp tránh data
leakage từ validation/test vào bộ biến đổi.

Trong M04, PCA giảm thời gian inference từ 5,93 ms/ảnh xuống 1,91 ms/ảnh
ở cấu hình validation, trong khi Macro F1 tăng từ 94,85\% lên 95,04\%.
Kết quả này cho thấy PCA vừa giảm chi phí tính toán vừa có tác dụng
regularization thực nghiệm trên bộ đặc trưng cụ thể.

\subsection{4.7. Tăng cường dữ liệu cho lớp
hiếm}\label{tux103ng-cux1b0ux1eddng-dux1eef-liux1ec7u-cho-lux1edbp-hiux1ebfm}

Đối với 10 lớp có dưới 30 ảnh train, nhóm áp dụng xoay nhẹ khoảng ±10°
và điều chỉnh độ sáng khoảng ±15\% trên ảnh 2D trước khi trích xuất đặc
trưng. Tổng số ảnh train tăng thêm 430. Cách làm này đảm bảo mỗi mẫu mới
tương ứng với một ảnh có ý nghĩa vật lý và HOG được tính lại từ cấu trúc
cạnh sau biến đổi.

SMOTE không được dùng trong pipeline cuối. Trong không gian HOG, nội suy
tuyến tính giữa hai vector có thể tạo mẫu không tương ứng với một ảnh cụ
thể. Điều này không có nghĩa SMOTE luôn không phù hợp, nhưng
augmentation ảnh là lựa chọn dễ kiểm soát và diễn giải hơn cho dự án.

\subsection{4.8. Bộ phân loại cuối
cùng}\label{bux1ed9-phuxe2n-loux1ea1i-cuux1ed1i-cuxf9ng}

Sau benchmark và tuning, mô hình được chọn là SVC(kernel="linear",
C=0.1, class\_weight=None). C=0,1 tạo regularization mạnh hơn các giá
trị C lớn, cho phép một số điểm vi phạm lề để giảm nhạy với outlier. Ba
chiến lược class\_weight None, balanced và custom cho cùng Validation
Macro F1 97,30\%; vì vậy None được chọn do đơn giản hơn và cho thấy
augmentation đã giảm nhu cầu điều chỉnh trọng số lớp trong split hiện
tại.

\section{5. THIẾT KẾ THỰC
NGHIỆM}\label{thiux1ebft-kux1ebf-thux1ef1c-nghiux1ec7m}

\subsection{5.1. Quy tắc fit và
transform}\label{quy-tux1eafc-fit-vuxe0-transform}

\begin{itemize}
\item
  Augmentation chỉ áp dụng trên train; validation và test không được
  tăng cường.
\item
  StandardScaler và PCA được fit trên train rồi mới transform
  validation/test.
\item
  Lựa chọn kernel và siêu tham số M06 dựa trên Validation Macro F1.
\item
  Macro F1 là metric chính; Accuracy và các metric còn lại dùng để giải
  thích trade-off.
\item
  Thời gian inference được chuẩn hóa theo ms/ảnh và cần ghi rõ phạm vi
  đo.
\end{itemize}

\subsection{5.2. Chỉ số đánh
giá}\label{chux1ec9-sux1ed1-ux111uxe1nh-giuxe1}

\begin{align}                                            
\text{Precision}_k &= \frac{\text{TP}_k}{\text{TP}_k +
\text{FP}_k} \\[0.6em]                                     
\text{Recall}_k &= \frac{\text{TP}_k}{\text{TP}_k +  
\text{FN}_k} \\[0.6em]                                     
\text{F1}_k &= \frac{2 \cdot \text{Precision}_k \cdot
\text{Recall}_k}{\text{Precision}_k + \text{Recall}_k}     
\\[0.8em]                                                  
\text{Macro F1} &= \frac{1}{K} \sum_{k=1}^{K}        
\text{F1}_k                                                
\end{align}

Macro F1 cho trọng số ngang nhau cho 47 lớp và được dùng làm tiêu chí
lựa chọn chính. Macro Recall đặc biệt hữu ích để quan sát mức bỏ sót các
lớp khó, còn Macro Precision phản ánh số false positive. Weighted F1 và
Accuracy giúp mô tả hiệu năng theo phân bố mẫu thực tế.

\subsection{5.3. Ma trận thực
nghiệm}\label{ma-trux1eadn-thux1ef1c-nghiux1ec7m}

\begin{longtable}[]{@{}
  >{\raggedright\arraybackslash}p{(\columnwidth - 4\tabcolsep) * \real{0.3333}}
  >{\raggedright\arraybackslash}p{(\columnwidth - 4\tabcolsep) * \real{0.3333}}
  >{\raggedright\arraybackslash}p{(\columnwidth - 4\tabcolsep) * \real{0.3333}}@{}}
\toprule\noalign{}
\begin{minipage}[b]{\linewidth}\raggedright
\textbf{Nhiệm vụ}
\end{minipage} & \begin{minipage}[b]{\linewidth}\raggedright
\textbf{Mục tiêu}
\end{minipage} & \begin{minipage}[b]{\linewidth}\raggedright
\textbf{Thành phần chính}
\end{minipage} \\
\midrule\noalign{}
\endhead
\bottomrule\noalign{}
\endlastfoot
M01 & Chẩn đoán dữ liệu và lớp hiếm & Phân bố lớp, bottleneck, nhóm
lỗi \\
M02 & Tối ưu kích thước và resize & 32/48/64; stretch/pad\_square \\
M03 & Khảo sát đặc trưng màu và fusion & HOG Gray, HOG YUV, HOG + HSV \\
M04 & Đánh giá PCA và augmentation & Không PCA; PCA; PCA +
augmentation \\
M05 & Benchmark mô hình & SVC, RF, LightGBM, LR, KNN, MLP \\
M06 & Fine-tuning SVM & Kernel, C, gamma, class\_weight \\
\end{longtable}

\begin{center}\small\textit{Bảng 3. Cấu trúc nhiệm vụ thực nghiệm của dự án.}\end{center}

\subsection{5.4. Audit việc sử dụng tập
test}\label{audit-viux1ec7c-sux1eed-dux1ee5ng-tux1eadp-test}

\begin{longtable}[]{@{}
  >{\raggedright\arraybackslash}p{(\columnwidth - 0\tabcolsep) * \real{1.0000}}@{}}
\toprule\noalign{}
\begin{minipage}[b]{\linewidth}\raggedright
\textbf{Lưu ý về tính hợp lệ của protocol}

Các artifact hiện có đã ghi nhận test metric ở Phase 1 và M05 trước khi
báo cáo kết quả M06. Vì vậy, không nên mô tả tập test như một blind
holdout chỉ được truy cập đúng một lần trong toàn bộ vòng đời dự án.
Việc chọn cấu hình M06 dựa trên validation giúp giảm rủi ro overfit trực
tiếp vào test, nhưng repeated observation vẫn là một đe dọa tính hợp lệ.
Phiên bản tiếp theo nên khóa test, dùng validation/cross-validation cho
toàn bộ lựa chọn và chỉ mở test sau khi pipeline được đóng băng.
\end{minipage} \\
\midrule\noalign{}
\endhead
\bottomrule\noalign{}
\endlastfoot
\end{longtable}

\subsection{5.5. Phạm vi của phép đo độ
trễ}\label{phux1ea1m-vi-cux1ee7a-phuxe9p-ux111o-ux111ux1ed9-trux1ec5}

Con số 1,52 ms/ảnh trong final result phản ánh thời gian dự đoán của
classifier trên vector đã được StandardScaler/PCA xử lý theo artifact
M06. Nó chưa bao gồm đọc ảnh, pad\_square, resize, HOG, HSV Histogram,
scaler và PCA. Do đó không nên quy đổi trực tiếp thành FPS end-to-end.
Muốn đánh giá triển khai thực tế, cần profiling từng bước trên cùng phần
cứng mục tiêu.

\section{6. KẾT QUẢ THỰC
NGHIỆM}\label{kux1ebft-quux1ea3-thux1ef1c-nghiux1ec7m}

\subsection{6.1. Kích thước ảnh và chế độ resize
(M02)}\label{kuxedch-thux1b0ux1edbc-ux1ea3nh-vuxe0-chux1ebf-ux111ux1ed9-resize-m02}

\begin{longtable}[]{@{}
  >{\raggedright\arraybackslash}p{(\columnwidth - 8\tabcolsep) * \real{0.2000}}
  >{\raggedright\arraybackslash}p{(\columnwidth - 8\tabcolsep) * \real{0.2000}}
  >{\raggedright\arraybackslash}p{(\columnwidth - 8\tabcolsep) * \real{0.2000}}
  >{\raggedright\arraybackslash}p{(\columnwidth - 8\tabcolsep) * \real{0.2000}}
  >{\raggedright\arraybackslash}p{(\columnwidth - 8\tabcolsep) * \real{0.2000}}@{}}
\toprule\noalign{}
\begin{minipage}[b]{\linewidth}\raggedright
\textbf{Cấu hình}
\end{minipage} & \begin{minipage}[b]{\linewidth}\raggedright
\textbf{Số chiều HOG}
\end{minipage} & \begin{minipage}[b]{\linewidth}\raggedright
\textbf{Val Accuracy (\%)}
\end{minipage} & \begin{minipage}[b]{\linewidth}\raggedright
\textbf{Val Macro F1 (\%)}
\end{minipage} & \begin{minipage}[b]{\linewidth}\raggedright
\textbf{Train (s)}
\end{minipage} \\
\midrule\noalign{}
\endhead
\bottomrule\noalign{}
\endlastfoot
32×32 stretch & 1.764 & 95,15 & 92,16 & 18,51 \\
32×32 pad\_square & 1.764 & 94,90 & 92,19 & 15,72 \\
48×48 stretch & 900 & 96,48 & 95,95 & 5,84 \\
48×48 pad\_square & 900 & 96,12 & 94,56 & 5,87 \\
64×64 stretch & 1.764 & 95,87 & 94,04 & 13,07 \\
64×64 pad\_square & 1.764 & 95,51 & 93,44 & 13,52 \\
\end{longtable}

\begin{center}\small\textit{Bảng 4. Kết quả khảo sát kích thước và chế độ resize ở M02 {[}5{]}.}\end{center}

48×48 stretch đạt Macro F1 cao nhất trong ablation HOG Gray. Kết quả cho
thấy tăng kích thước không tự động làm hiệu năng tăng; số chiều lớn hơn
có thể thêm nhiễu và tăng chi phí. Tuy nhiên, bảng này chưa đánh giá
toàn bộ pipeline có màu, PCA và augmentation nên không trực tiếp quyết
định cấu hình final.

\subsection{6.2. Không gian màu và dung hợp đặc trưng
(M03)}\label{khuxf4ng-gian-muxe0u-vuxe0-dung-hux1ee3p-ux111ux1eb7c-trux1b0ng-m03}

\begin{longtable}[]{@{}
  >{\raggedright\arraybackslash}p{(\columnwidth - 8\tabcolsep) * \real{0.2000}}
  >{\raggedright\arraybackslash}p{(\columnwidth - 8\tabcolsep) * \real{0.2000}}
  >{\raggedright\arraybackslash}p{(\columnwidth - 8\tabcolsep) * \real{0.2000}}
  >{\raggedright\arraybackslash}p{(\columnwidth - 8\tabcolsep) * \real{0.2000}}
  >{\raggedright\arraybackslash}p{(\columnwidth - 8\tabcolsep) * \real{0.2000}}@{}}
\toprule\noalign{}
\begin{minipage}[b]{\linewidth}\raggedright
\textbf{Đặc trưng}
\end{minipage} & \begin{minipage}[b]{\linewidth}\raggedright
\textbf{Số chiều}
\end{minipage} & \begin{minipage}[b]{\linewidth}\raggedright
\textbf{Val Accuracy (\%)}
\end{minipage} & \begin{minipage}[b]{\linewidth}\raggedright
\textbf{Val Macro F1 (\%)}
\end{minipage} & \begin{minipage}[b]{\linewidth}\raggedright
\textbf{Train (s)}
\end{minipage} \\
\midrule\noalign{}
\endhead
\bottomrule\noalign{}
\endlastfoot
HOG Gray & 900 & 96,48 & 95,95 & 6,33 \\
HOG YUV & 2.700 & 96,36 & 94,75 & 25,60 \\
HOG Gray + HSV Histogram & 1.412 & 95,87 & 95,19 & 9,08 \\
\end{longtable}

\begin{center}\small\textit{Bảng 5. Kết quả khảo sát đặc trưng ở M03 {[}5{]}.}\end{center}

HOG Gray đơn lẻ đạt F1 cao nhất trong ablation M03, còn HOG YUV vừa tăng
số chiều vừa giảm F1. HOG + HSV chưa vượt HOG Gray ở cấu hình 48×48
riêng lẻ nhưng được giữ lại vì cung cấp thông tin màu bổ sung và cho kết
quả tốt hơn khi kết hợp với cấu hình M04, PCA và augmentation. Điều này
cho thấy không nên chọn thành phần chỉ dựa trên một thí nghiệm cô lập.

\subsection{6.3. Tác động của PCA và augmentation
(M04)}\label{tuxe1c-ux111ux1ed9ng-cux1ee7a-pca-vuxe0-augmentation-m04}

\begin{longtable}[]{@{}
  >{\raggedright\arraybackslash}p{(\columnwidth - 10\tabcolsep) * \real{0.1667}}
  >{\raggedright\arraybackslash}p{(\columnwidth - 10\tabcolsep) * \real{0.1667}}
  >{\raggedright\arraybackslash}p{(\columnwidth - 10\tabcolsep) * \real{0.1667}}
  >{\raggedright\arraybackslash}p{(\columnwidth - 10\tabcolsep) * \real{0.1667}}
  >{\raggedright\arraybackslash}p{(\columnwidth - 10\tabcolsep) * \real{0.1667}}
  >{\raggedright\arraybackslash}p{(\columnwidth - 10\tabcolsep) * \real{0.1667}}@{}}
\toprule\noalign{}
\begin{minipage}[b]{\linewidth}\raggedright
\textbf{Cấu hình}
\end{minipage} & \begin{minipage}[b]{\linewidth}\raggedright
\textbf{Số chiều}
\end{minipage} & \begin{minipage}[b]{\linewidth}\raggedright
\textbf{Số ảnh train}
\end{minipage} & \begin{minipage}[b]{\linewidth}\raggedright
\textbf{Val Accuracy (\%)}
\end{minipage} & \begin{minipage}[b]{\linewidth}\raggedright
\textbf{Val Macro F1 (\%)}
\end{minipage} & \begin{minipage}[b]{\linewidth}\raggedright
\textbf{Inference (ms/ảnh)}
\end{minipage} \\
\midrule\noalign{}
\endhead
\bottomrule\noalign{}
\endlastfoot
TN3.1: Không PCA, không Aug & 1.812 & 6.605 & 96,24 & 94,85 & 5,93 \\
TN3.2: PCA 95\% & 569 & 6.605 & 96,60 & 95,04 & 1,91 \\
TN3.3: PCA 95\% + Aug & 571 & 7.035 & 96,72 & 95,77 & 1,96 \\
\end{longtable}

\begin{center}\small\textit{Bảng 6. Tác động của PCA và minority augmentation ở M04 {[}5{]}.}\end{center}

PCA giảm khoảng 68,5\% số chiều và làm inference nhanh hơn hơn ba lần
trong ablation M04. Augmentation bổ sung 430 ảnh và nâng Macro F1
validation từ 95,04\% lên 95,77\%. Artifact Phase 1 ghi nhận mô hình RBF
đạt Test Accuracy 96,24\%, Test Macro F1 95,64\% và 1,93 ms/ảnh.

\subsection{6.4. Benchmark sáu họ mô hình
(M05)}\label{benchmark-suxe1u-hux1ecd-muxf4-huxecnh-m05}

\begin{longtable}[]{@{}
  >{\raggedright\arraybackslash}p{(\columnwidth - 10\tabcolsep) * \real{0.1667}}
  >{\raggedright\arraybackslash}p{(\columnwidth - 10\tabcolsep) * \real{0.1667}}
  >{\raggedright\arraybackslash}p{(\columnwidth - 10\tabcolsep) * \real{0.1667}}
  >{\raggedright\arraybackslash}p{(\columnwidth - 10\tabcolsep) * \real{0.1667}}
  >{\raggedright\arraybackslash}p{(\columnwidth - 10\tabcolsep) * \real{0.1667}}
  >{\raggedright\arraybackslash}p{(\columnwidth - 10\tabcolsep) * \real{0.1667}}@{}}
\toprule\noalign{}
\begin{minipage}[b]{\linewidth}\raggedright
\textbf{Mô hình}
\end{minipage} & \begin{minipage}[b]{\linewidth}\raggedright
\textbf{Train (s)}
\end{minipage} & \begin{minipage}[b]{\linewidth}\raggedright
\textbf{Inference (ms/ảnh)}
\end{minipage} & \begin{minipage}[b]{\linewidth}\raggedright
\textbf{Accuracy (\%)}
\end{minipage} & \begin{minipage}[b]{\linewidth}\raggedright
\textbf{Macro F1 (\%)}
\end{minipage} & \begin{minipage}[b]{\linewidth}\raggedright
\textbf{Minority Recall (\%)}
\end{minipage} \\
\midrule\noalign{}
\endhead
\bottomrule\noalign{}
\endlastfoot
SVC RBF & 3,5868 & 2,4520 & 95,77 & 95,37 & 86,67 \\
Random Forest & 2,6118 & 0,0310 & 85,78 & 84,74 & 75,00 \\
LightGBM & 40,4803 & 0,0373 & 87,78 & 86,70 & 80,83 \\
Logistic Regression & 1,2973 & 0,0027 & 95,77 & 94,82 & 89,17 \\
KNN & 0,0117 & 1,8068 & 94,83 & 94,23 & 87,50 \\
MLP & 2,9231 & 0,0018 & 95,18 & 93,62 & 87,50 \\
\end{longtable}

\begin{center}\small\textit{Bảng 7. Benchmark sáu mô hình trên đặc trưng PCA 571 chiều {[}6{]}.}\end{center}

\begin{figure}[H]
\centering
\includegraphics[width=0.92\textwidth]{./media/image4.png}
\caption*{\textit{Hình 4. So sánh Macro F1 của sáu họ mô hình trong artifact M05.}}
\end{figure}

SVC RBF dẫn đầu benchmark; Logistic Regression đạt Accuracy tương đương
và Macro F1 chỉ thấp hơn 0,55 điểm phần trăm. Đây là tín hiệu rằng không
gian sau PCA phù hợp tốt với mô hình tuyến tính. Random Forest và
LightGBM có inference nhanh nhưng F1 thấp hơn đáng kể; KNN và MLP cạnh
tranh nhưng không vượt SVM.

\subsection{6.5. Khảo sát kernel và siêu tham số SVM
(M06)}\label{khux1ea3o-suxe1t-kernel-vuxe0-siuxeau-tham-sux1ed1-svm-m06}

\begin{longtable}[]{@{}
  >{\raggedright\arraybackslash}p{(\columnwidth - 8\tabcolsep) * \real{0.2000}}
  >{\raggedright\arraybackslash}p{(\columnwidth - 8\tabcolsep) * \real{0.2000}}
  >{\raggedright\arraybackslash}p{(\columnwidth - 8\tabcolsep) * \real{0.2000}}
  >{\raggedright\arraybackslash}p{(\columnwidth - 8\tabcolsep) * \real{0.2000}}
  >{\raggedright\arraybackslash}p{(\columnwidth - 8\tabcolsep) * \real{0.2000}}@{}}
\toprule\noalign{}
\begin{minipage}[b]{\linewidth}\raggedright
\textbf{Kernel / cấu hình}
\end{minipage} & \begin{minipage}[b]{\linewidth}\raggedright
\textbf{Val Accuracy (\%)}
\end{minipage} & \begin{minipage}[b]{\linewidth}\raggedright
\textbf{Val Macro F1 (\%)}
\end{minipage} & \begin{minipage}[b]{\linewidth}\raggedright
\textbf{Train (s)}
\end{minipage} & \begin{minipage}[b]{\linewidth}\raggedright
\textbf{Support vectors}
\end{minipage} \\
\midrule\noalign{}
\endhead
\bottomrule\noalign{}
\endlastfoot
Linear, C=0,1,\newline class\_weight=\newline None & 96,72 & 97,30 & 2,46 & 4.514 \\
Linear, C=10, balanced & 96,72 & 97,30 & 2,22 & 4.514 \\
RBF, C=10, gamma=scale & 96,72 & 95,77 & 6,25 & 5.331 \\
RBF, C=5,\newline gamma=\newline 0,0005 & 96,84 & 95,81 & 5,20 & 5.261 \\
Polynomial degree 2 & 96,72 & 95,03 & 6,32 & 5.273 \\
Polynomial degree 3 & 94,78 & 93,43 & 14,28 & 5.172 \\
Sigmoid & 94,05 & 92,96 & 2,35 & 3.812 \\
\end{longtable}

\begin{center}\small\textit{Bảng 8. Kết quả khảo sát kernel SVM trên validation {[}6{]}, {[}7{]}.}\end{center}

\begin{figure}[H]
\centering
\includegraphics[width=0.92\textwidth]{./media/image5.png}
\caption*{\textit{Hình 5. Validation Macro F1 theo kernel SVM.}}
\end{figure}

Linear SVM đạt Validation Macro F1 97,30\%, cao hơn RBF baseline 1,53
điểm phần trăm. Nhiều giá trị C cho cùng metric validation, cho thấy
nghiệm khá ổn định trên split này. C=0,1 được chọn theo nguyên tắc ưu
tiên regularization mạnh hơn khi hiệu năng tương đương. Ba chiến lược
class\_weight cũng cho cùng F1, vì vậy class\_weight=None được chọn.

\subsection{6.6. Kết quả chung cuộc và
trade-off}\label{kux1ebft-quux1ea3-chung-cuux1ed9c-vuxe0-trade-off}

\begin{longtable}[]{@{}
  >{\raggedright\arraybackslash}p{(\columnwidth - 6\tabcolsep) * \real{0.2500}}
  >{\raggedright\arraybackslash}p{(\columnwidth - 6\tabcolsep) * \real{0.2500}}
  >{\raggedright\arraybackslash}p{(\columnwidth - 6\tabcolsep) * \real{0.2500}}
  >{\raggedright\arraybackslash}p{(\columnwidth - 6\tabcolsep) * \real{0.2500}}@{}}
\toprule\noalign{}
\begin{minipage}[b]{\linewidth}\raggedright
\textbf{Chỉ số}
\end{minipage} & \begin{minipage}[b]{\linewidth}\raggedright
\textbf{Phase 1: RBF}
\end{minipage} & \begin{minipage}[b]{\linewidth}\raggedright
\textbf{Phase 2: Linear}
\end{minipage} & \begin{minipage}[b]{\linewidth}\raggedright
\textbf{Chênh lệch}
\end{minipage} \\
\midrule\noalign{}
\endhead
\bottomrule\noalign{}
\endlastfoot
Test Accuracy (\%) & 96,24 & 96,12 & -0,12 \\
Test Macro F1 (\%) & 95,64 & 96,29 & +0,65 \\
Test Macro Precision (\%) & 98,08 & 97,51 & -0,57 \\
Test Macro Recall (\%) & 94,39 & 95,62 & +1,23 \\
Test Weighted F1 (\%) & 96,15 & 96,09 & -0,06 \\
Classifier inference (ms/ảnh) & 1,93 & 1,52 & -0,41 \\
Support vectors & 5.331 & 4.514 & -817 \\
\end{longtable}

\begin{center}\small\textit{Bảng 9. So sánh kết quả được ghi nhận giữa Phase 1 và Phase 2 {[}6{]}.}\end{center}

\begin{figure}[H]
\centering
\includegraphics[width=0.98\textwidth]{./media/image8.png}
\caption*{\textit{Hình 6. Diễn biến cải thiện kết quả qua bốn mốc phát triển của hệ thống.}}
\end{figure}

Biểu đồ ở Hình 6 cho thấy hiệu năng được cải thiện qua từng giai đoạn. Từ baseline ban đầu, PCA vừa giảm mạnh số chiều từ
1.812 xuống 569 vừa tăng nhẹ Validation Macro F1. Sau đó, minority
augmentation tiếp tục nâng Validation Macro F1 lên 95,77\% và giúp cải
thiện khả năng nhận diện lớp hiếm. Ở giai đoạn cuối, việc chuyển từ RBF
SVM sang Linear SVM trên không gian đặc trưng đã được chuẩn hóa và giảm
chiều tạo ra bước nhảy lớn nhất, đưa Validation Macro F1 lên 97,30\% và
Test Macro F1 lên 96,29\%.

\begin{figure}[H]
\centering
\includegraphics[width=0.92\textwidth]{./media/image6.png}
\caption*{\textit{Hình 7. Đối chiếu các metric giữa Phase 1 và Phase 2.}}
\end{figure}

Linear SVM tăng Macro F1 0,65 điểm phần trăm và Macro Recall 1,23 điểm
phần trăm, trong khi Accuracy giảm 0,12 và Macro Precision giảm 0,57.
Đây là trade-off hợp lý trong dữ liệu mất cân bằng: mô hình nhận diện
tốt hơn các lớp khó và giảm false negative, đổi lại phát sinh thêm một
số false positive. Vì Macro F1 là metric chính, Linear SVM được chọn làm
mô hình cuối.

\subsection{6.7. Phân tích nhóm lỗi còn
lại}\label{phuxe2n-tuxedch-nhuxf3m-lux1ed7i-cuxf2n-lux1ea1i}

\begin{longtable}[]{@{}
  >{\raggedright\arraybackslash}p{(\columnwidth - 6\tabcolsep) * \real{0.2500}}
  >{\raggedright\arraybackslash}p{(\columnwidth - 6\tabcolsep) * \real{0.2500}}
  >{\raggedright\arraybackslash}p{(\columnwidth - 6\tabcolsep) * \real{0.2500}}
  >{\raggedright\arraybackslash}p{(\columnwidth - 6\tabcolsep) * \real{0.2500}}@{}}
\toprule\noalign{}
\begin{minipage}[b]{\linewidth}\raggedright
\textbf{Cặp / nhóm lớp}
\end{minipage} & \begin{minipage}[b]{\linewidth}\raggedright
\textbf{Đặc điểm gây nhầm}
\end{minipage} & \begin{minipage}[b]{\linewidth}\raggedright
\textbf{Nguyên nhân khả dĩ}
\end{minipage} & \begin{minipage}[b]{\linewidth}\raggedright
\textbf{Hướng xử lý}
\end{minipage} \\
\midrule\noalign{}
\endhead
\bottomrule\noalign{}
\endlastfoot
P.127\_60 $\leftrightarrow$ P.127\_80 & Cùng biển tròn viền đỏ, khác chữ số & Chi tiết
số nhỏ, mờ hoặc nén bởi PCA & Crop vùng số; đặc trưng độ phân giải cao;
classifier chuyên biệt \\
W.224 $\leftrightarrow$ W.245a & Cùng tam giác nền vàng, biểu tượng đen & Hình dạng tổng
thể giống; biểu tượng trung tâm nhỏ & Đặc trưng local; multi-stage
classification \\
P.130 $\leftrightarrow$ P.131a & Cấu trúc hình học gần giống, khác hướng ký hiệu & Lệch
góc và blur làm mất hướng & Augmentation góc nhìn; orientation-sensitive
local features \\
W.205c & Lớp cực hiếm & 17 ảnh train, biến thiên test chưa được đại diện
& Thu thập dữ liệu thật; few-shot/metric learning trong tương lai \\
\end{longtable}

\begin{center}\small\textit{Bảng 10. Các nhóm lỗi tiêu biểu và hướng khắc phục.}\end{center}

Các lỗi trên cho thấy giới hạn không nằm hoàn toàn ở bộ phân loại. Khi
hai lớp có cùng hình dạng và màu sắc, phần thông tin quyết định thường
chỉ nằm ở một vùng rất nhỏ. Nếu vùng này bị mờ, nén hoặc chiếm quá ít
pixel, cả HOG và histogram màu đều khó cung cấp tín hiệu đủ mạnh. Vì
vậy, tăng độ phức tạp của classifier chưa chắc giải quyết được nguyên
nhân gốc.

W.205c minh họa rõ small-sample problem: chỉ 17 ảnh train không đủ bao
phủ các góc nhìn và điều kiện ảnh khác nhau. Hướng ưu tiên nên là bổ
sung dữ liệu thật và đánh giá per-class ổn định hơn, trước khi áp dụng
các mô hình phức tạp hoặc kỹ thuật tổng hợp mạnh.

\section{7. THẢO LUẬN}\label{thux1ea3o-luux1eadn}

\subsection{7.1. Vì sao HOG và HSV bổ sung cho
nhau?}\label{vuxec-sao-hog-vuxe0-hsv-bux1ed5-sung-cho-nhau}

HOG tập trung vào hình dạng, đường viền và biểu tượng; HSV Histogram mô
tả phân bố màu tổng thể. Hai biển có hình dạng giống nhau nhưng màu khác
có thể được tách nhờ HSV, trong khi hai biển cùng màu nền nhưng biểu
tượng khác nhau được tách nhờ HOG. Sự bổ sung này phù hợp với cấu trúc
thiết kế của biển báo giao thông. Tuy nhiên, ablation M03 cho thấy
fusion không tự động tốt hơn HOG đơn lẻ; lợi ích chỉ xuất hiện khi xét
cả pipeline cùng PCA và augmentation.

\subsection{7.2. Vai trò thực nghiệm của
PCA}\label{vai-truxf2-thux1ef1c-nghiux1ec7m-cux1ee7a-pca}

PCA giảm bộ nhớ, số phép tính và độ trễ, đồng thời cải thiện nhẹ Macro
F1 trong M04. Tuy nhiên, PCA tối đa hóa phương sai chứ không tối ưu trực
tiếp khả năng phân lớp. Một số chi tiết có phương sai nhỏ nhưng quan
trọng cho nhãn, ví dụ nét chữ số, có thể bị suy giảm khi nén. Vì vậy,
kết luận đúng là ``PCA hiệu quả trên bộ đặc trưng và split này'', không
phải ``PCA luôn loại bỏ nhiễu và luôn tăng độ chính xác''.

\subsection{7.3. Vì sao Linear SVM tốt hơn RBF trong pipeline
này?}\label{vuxec-sao-linear-svm-tux1ed1t-hux1a1n-rbf-trong-pipeline-nuxe0y}

Kết quả Logistic Regression và Linear SVM cho thấy dữ liệu sau
StandardScaler và PCA có khả năng phân tách tốt bằng ranh giới tuyến
tính. Khi đó, RBF thêm độ linh hoạt không cần thiết và có thể nhạy hơn
với tham số. C=0,1 tạo soft margin rộng hơn, giảm nhạy với outlier và
ảnh biến dạng. Kết luận chỉ áp dụng cho pipeline HOG+HSV+PCA hiện tại;
không thể suy rộng rằng Linear SVM luôn tốt hơn RBF trong mọi bài toán
biển báo.

Về triển khai, SVC kernel tuyến tính vẫn có thể dựa trên support vectors
trong quá trình dự đoán tùy cách hiện thực của thư viện. Vì vậy báo cáo
sử dụng thời gian đo thực tế và số support vector, không đơn giản hóa
thành giả định rằng dự đoán chắc chắn chỉ là một phép $w^{\mathsf{T}}x+b$ duy nhất
trong mọi implementation.

\subsection{7.4. Vì sao Random Forest và LightGBM thấp
hơn?}\label{vuxec-sao-random-forest-vuxe0-lightgbm-thux1ea5p-hux1a1n}

Cây quyết định chia không gian bằng các ngưỡng theo từng chiều. Sau PCA,
mỗi chiều là tổ hợp tuyến tính của nhiều đặc trưng gốc; ranh giới lớp có
thể nằm theo hướng chéo và cần nhiều phép chia để xấp xỉ, tạo vùng quyết
định phân mảnh. Đây là một giải thích hợp lý cho việc Random Forest và
LightGBM thấp hơn SVM 8,7-10,6 điểm Macro F1, nhưng nên được xem là giả
thuyết dựa trên đặc trưng của mô hình và kết quả, không phải chứng minh
duy nhất.

\subsection{7.5. Hạn chế của KNN trong 571
chiều}\label{hux1ea1n-chux1ebf-cux1ee7a-knn-trong-571-chiux1ec1u}

KNN dựa vào khoảng cách giữa mẫu. Trong không gian cao chiều, khoảng
cách có xu hướng ít phân biệt hơn và dự đoán phải so sánh với nhiều mẫu
train. PCA đã giảm chiều đáng kể nhưng 571 vẫn cao. KNN đạt Macro F1
94,23\%, cho thấy đặc trưng có chất lượng, nhưng không vượt SVM và
classifier inference khoảng 1,81 ms/ảnh.

\subsection{7.6. Ý nghĩa của việc tăng Macro
Recall}\label{uxfd-nghux129a-cux1ee7a-viux1ec7c-tux103ng-macro-recall}

Phase 2 giảm nhẹ Accuracy và Precision nhưng tăng Recall. Trong ứng dụng
biển báo, bỏ sót một biển hiếm có thể nghiêm trọng hơn một số false
positive bổ sung, tùy kịch bản triển khai. Macro Recall tăng 1,23 điểm
cho thấy mô hình quan tâm tốt hơn đến các lớp khó. Tuy nhiên, để kết
luận về an toàn, cần đánh giá theo chi phí lỗi của từng lớp và trên dữ
liệu đường phố thực tế, không chỉ dựa vào metric trung bình.

\subsection{7.7. Khả năng triển
khai}\label{khux1ea3-nux103ng-triux1ec3n-khai}

Classifier inference 1,52 ms/ảnh tương đương khoảng 658 dự đoán/giây nếu
chỉ xét bước phân loại. Đây không phải FPS end-to-end. Dù vậy, kết quả
cho thấy classifier không phải nút thắt lớn và pipeline có tiềm năng
thời gian thực nếu HOG, HSV và PCA được tối ưu. Cần đo trên phần cứng
mục tiêu, báo cáo RAM, kích thước model và độ trễ từng stage trước khi
khẳng định khả năng triển khai edge.

\section{8. HẠN CHẾ VÀ ĐE DỌA TÍNH HỢP
LỆ}\label{hux1ea1n-chux1ebf-vuxe0-ux111e-dux1ecda-tuxednh-hux1ee3p-lux1ec7}

\subsection{8.1. Hạn chế dữ
liệu}\label{hux1ea1n-chux1ebf-dux1eef-liux1ec7u}

\begin{itemize}
\item
  Một số lớp có quá ít ảnh, khiến metric per-class có phương sai cao và
  khó khái quát.
\item
  Ảnh đã crop nên chưa đánh giá sai số detection, background clutter và
  multiple signs.
\item
  Augmentation mới bao phủ xoay và sáng nhẹ; chưa bao phủ perspective,
  motion blur, occlusion và weather.
\item
  Một split duy nhất chưa đủ để ước lượng độ ổn định của mô hình.
\end{itemize}

\subsection{8.2. Hạn chế protocol}\label{hux1ea1n-chux1ebf-protocol}

\begin{itemize}
\item
  Các artifact đã quan sát test ở nhiều giai đoạn; test không còn là
  blind holdout tuyệt đối.
\item
  Một số benchmark M05 dùng cột tên ``Test'', vì vậy cần tránh sử dụng
  chính các kết quả đó để lựa chọn lặp lại mà không có validation độc
  lập.
\item
  Chưa báo cáo khoảng tin cậy hoặc kiểm định thống kê cho chênh lệch
  0,65 điểm Macro F1.
\item
  Thời gian train/inference có thể không so sánh trực tiếp nếu các bước
  đo hoặc phần cứng khác nhau.
\end{itemize}

\subsection{8.3. Hạn chế tái lập và nhất quán
artifact}\label{hux1ea1n-chux1ebf-tuxe1i-lux1eadp-vuxe0-nhux1ea5t-quuxe1n-artifact}

\begin{itemize}
\item
  Số chiều 1.812 được lấy từ artifact M04 nhưng chưa khớp hoàn toàn với
  decomposition lý thuyết từ các tham số HOG được ghi lại.
\item
  Cần lưu phiên bản thư viện, seed, phần cứng, mã hash dữ liệu và
  command chạy cuối để tái lập đầy đủ.
\item
  Con số 1,52 ms/ảnh chỉ là classifier inference; chưa có profiling
  end-to-end.
\item
  Chưa báo cáo kích thước file model, peak RAM và thời gian khởi tạo
  pipeline.
\end{itemize}

\section{9. HƯỚNG PHÁT TRIỂN}\label{hux1b0ux1edbng-phuxe1t-triux1ec3n}

\begin{enumerate}
\def\labelenumi{\arabic{enumi}.}
\setcounter{enumi}{10}
\item
  Thu thập thêm dữ liệu thật cho các lớp dưới 30 ảnh, ưu tiên W.205c và
  các lớp có per-class Recall thấp.
\item
  Xây dựng phân loại phân cấp coarse-to-fine: tầng đầu nhận diện nhóm
  hình dạng/màu, tầng sau xử lý các cặp dễ nhầm, tham khảo chiến lược
  của Ma và Huang {[}3{]}.
\item
  Tách vùng trung tâm hoặc vùng chữ số ở độ phân giải cao cho nhóm biển
  tốc độ.
\item
  Mở rộng augmentation có kiểm soát với perspective transform, blur,
  contrast, occlusion và mô phỏng thời tiết.
\item
  Thực hiện repeated stratified cross-validation hoặc nested validation
  cho toàn pipeline; khóa test chỉ dùng một lần.
\item
  So sánh SVC linear với LinearSVC, SGDClassifier hoặc mô hình tuyến
  tính nén để giảm kích thước và latency.
\item
  Đo end-to-end latency, RAM và model size trên CPU/edge device, phân rã
  thời gian từng stage.
\item
  Kết hợp module phát hiện biển báo để tạo pipeline end-to-end trên ảnh
  đường phố.
\end{enumerate}

\section{10. KẾT LUẬN}\label{kux1ebft-luux1eadn}

Dự án đã xây dựng một pipeline nhận diện biển báo giao thông Việt Nam
dựa trên computer vision truyền thống và học máy cổ điển. Pipeline cuối
gồm 64×64 pad\_square, HOG Gray, HSV Color Histogram, StandardScaler,
PCA giữ 95\% phương sai, augmentation ảnh cho 10 lớp thiểu số và SVC
tuyến tính với C=0,1.

Quá trình thực nghiệm được tổ chức theo hai giai đoạn. Phase 1 tối ưu
tiền xử lý và đặc trưng với RBF cố định; Phase 2 benchmark sáu họ mô
hình và tinh chỉnh SVM. Cấu hình cuối đạt Accuracy 96,12\%, Macro F1
96,29\%, Macro Recall 95,62\% và classifier inference 1,52 ms/ảnh trên
tập test 851 ảnh. So với RBF Phase 1, Linear SVM tăng Macro F1 0,65 điểm
phần trăm và Macro Recall 1,23 điểm phần trăm, đồng thời giảm 817
support vector.

Kết quả chứng minh rằng thiết kế biểu diễn có thể quan trọng hơn việc
tăng độ phức tạp mô hình. HOG và HSV cung cấp thông tin bổ sung; PCA
giảm mạnh số chiều; augmentation giúp lớp hiếm; và Linear SVM phù hợp
tốt với không gian sau biến đổi. Đồng thời, báo cáo xác định rõ các giới
hạn về repeated test exposure, latency chưa end-to-end, lớp cực hiếm và
sự chưa thống nhất số chiều artifact. Những điểm này tạo nền tảng minh
bạch cho việc cải tiến và tái lập dự án trong giai đoạn tiếp theo.

\begin{longtable}[]{@{}
  >{\raggedright\arraybackslash}p{(\columnwidth - 0\tabcolsep) * \real{1.0000}}@{}}
\toprule\noalign{}
\begin{minipage}[b]{\linewidth}\raggedright
\textbf{Kết luận lựa chọn mô hình}

Linear SVM được chọn không phải vì mọi chỉ số đều cao hơn RBF, mà vì nó
tối ưu tiêu chí chính Macro F1, tăng Macro Recall, giảm số support
vector và có classifier inference thấp hơn. Đây là một quyết định dựa
trên trade-off, không phải chỉ dựa trên Accuracy.
\end{minipage} \\
\midrule\noalign{}
\endhead
\bottomrule\noalign{}
\endlastfoot
\end{longtable}

\section{TÀI LIỆU THAM KHẢO}\label{tuxe0i-liux1ec7u-tham-khux1ea3o}

{[}1{]} C. Cortes and V. Vapnik, ``Support-vector networks,'' Machine
Learning, vol. 20, pp. 273-297, 1995.

{[}2{]} N. Dalal and B. Triggs, ``Histograms of oriented gradients for
human detection,'' in Proc. IEEE CVPR, 2005, pp. 886-893.

{[}3{]} Y. Ma and L. Huang, ``Hierarchical traffic sign recognition
based on multi-feature and multi-classifier fusion,'' First
International Conference on Information Science and Electronic
Technology, 2015, pp. 56-59.

{[}4{]} J. Stallkamp, M. Schlipsing, J. Salmen, and C. Igel, ``Man vs.
computer: Benchmarking machine learning algorithms for traffic sign
recognition,'' Neural Networks, vol. 32, pp. 323-332, 2012.

{[}5{]} R01, ``Báo cáo tổng kết Phase 1-New: tối ưu hóa tiền xử lý và kỹ
thuật đặc trưng,'' tài liệu nội bộ dự án, 05/07/2026.

{[}6{]} R02, ``Báo cáo tổng kết Phase 2-New: khảo sát mô hình học máy và
tinh chỉnh SVM,'' tài liệu nội bộ dự án, 12/07/2026.

{[}7{]} M07, ``Tài liệu tổng hợp và bàn giao cho Phase 3 Paper
Writing,'' tài liệu nội bộ dự án, 12/07/2026.

\section{PHỤ LỤC A. CẤU HÌNH CUỐI CÙNG CỦA HỆ
THỐNG}\label{phux1ee5-lux1ee5c-a.-cux1ea5u-huxecnh-cuux1ed1i-cuxf9ng-cux1ee7a-hux1ec7-thux1ed1ng}

\begin{longtable}[]{@{}
  >{\raggedright\arraybackslash}p{(\columnwidth - 2\tabcolsep) * \real{0.5000}}
  >{\raggedright\arraybackslash}p{(\columnwidth - 2\tabcolsep) * \real{0.5000}}@{}}
\toprule\noalign{}
\begin{minipage}[b]{\linewidth}\raggedright
\textbf{Thành phần}
\end{minipage} & \begin{minipage}[b]{\linewidth}\raggedright
\textbf{Cấu hình chốt}
\end{minipage} \\
\midrule\noalign{}
\endhead
\bottomrule\noalign{}
\endlastfoot
Dữ liệu & NVTS, 47 lớp; Train 6.605, Validation 824, Test 851 \\
Tiền xử lý & 64×64, pad\_square \\
Đặc trưng hình dạng & HOG trên ảnh grayscale \\
Đặc trưng màu & HSV Color Histogram 8×8×8 bins \\
Dung hợp & Vector 1.812 chiều theo artifact M04 \\
Chuẩn hóa & StandardScaler fit trên train \\
Giảm chiều & PCA 95\% phương sai, 571 chiều final \\
Xử lý mất cân bằng & Augmentation ảnh 2D cho 10 lớp \textless30 mẫu;
train tăng lên 7.035 \\
Bộ phân loại & SVC kernel=linear, C=0,1, class\_weight=None \\
Metric chính & Macro F1 \\
Final result & Accuracy 96,12\%; Macro F1 96,29\%; Macro Recall
95,62\% \\
Inference & 1,52 ms/ảnh ở bước classifier trên đặc trưng đã xử lý \\
\end{longtable}

\begin{center}\small\textit{Bảng A1. Bộ tham số cuối cùng được sử dụng trong báo cáo.}\end{center}

\section{PHỤ LỤC B. BẢNG KIỂM SOÁT TÍNH NHẤT
QUÁN}\label{phux1ee5-lux1ee5c-b.-bux1ea3ng-kiux1ec3m-souxe1t-tuxednh-nhux1ea5t-quuxe1n}

\begin{longtable}[]{@{}
  >{\raggedright\arraybackslash}p{(\columnwidth - 4\tabcolsep) * \real{0.3333}}
  >{\raggedright\arraybackslash}p{(\columnwidth - 4\tabcolsep) * \real{0.3333}}
  >{\raggedright\arraybackslash}p{(\columnwidth - 4\tabcolsep) * \real{0.3333}}@{}}
\toprule\noalign{}
\begin{minipage}[b]{\linewidth}\raggedright
\textbf{Mục kiểm tra}
\end{minipage} & \begin{minipage}[b]{\linewidth}\raggedright
\textbf{Giá trị chuẩn}
\end{minipage} & \begin{minipage}[b]{\linewidth}\raggedright
\textbf{Ghi chú}
\end{minipage} \\
\midrule\noalign{}
\endhead
\bottomrule\noalign{}
\endlastfoot
Final Macro F1 & 96,29\% & Dùng thống nhất trong Tóm tắt, Kết quả và Kết
luận \\
Final Accuracy & 96,12\% & Không nhầm với 96,24\% của Phase 1 \\
Final Macro Recall & 95,62\% & Cải thiện +1,23 điểm phần trăm so với
Phase 1 \\
Classifier inference & 1,52 ms/ảnh & Không ghi là end-to-end latency \\
PCA dimension & 571 & Vector final sau PCA ở pipeline cuối \\
Feature dimension & 1.812 & Theo artifact M04; cần audit code nếu công
bố decomposition \\
Lớp hiếm & 10 lớp \textless30 ảnh & Không khẳng định tất cả đều đạt
100\% \\
Final model & SVC linear, C=0,1, class\_weight=None & Cấu hình chốt
M06 \\
Test protocol & Đã quan sát ở nhiều artifact & Không mô tả là blind
one-shot tuyệt đối \\
\end{longtable}

\begin{center}\small\textit{Bảng B1. Các giá trị cần giữ nhất quán khi chỉnh sửa hoặc trình bày.}\end{center}

\section{PHỤ LỤC C. ARTIFACT NỘI BỘ PHỤC VỤ TÁI
LẬP}\label{phux1ee5-lux1ee5c-c.-artifact-nux1ed9i-bux1ed9-phux1ee5c-vux1ee5-tuxe1i-lux1eadp}

Các đường dẫn dưới đây là artifact nội bộ của repository, không thay thế
tài liệu tham khảo học thuật nhưng cần được lưu để truy vết số liệu:

\begin{itemize}
\item
  \path{paper/result/report_new/R01_26.07.05_report-phase1-new.md}
\item
  \path{paper/result/report_new/R02_26.07.12_report-phase2-new.md}
\item
  \path{paper/mission/phase3_new/M07_26-07-12_master_handover_document_for_paper_writing.md}
\item
  \path{paper/workspace/TruongDT/phase2_new/outputs/M05_ml_models_comparison_results.csv}
\item
  \path{paper/workspace/KietBA/phase2_new/outputs/M06_svm_finetuning_results.csv}
\item
  \path{paper/workspace/KietBA/phase2_new/outputs/M06_final_test_summary.csv}
\end{itemize}

\subsection{C.1. Checklist tái lập trước khi
nộp}\label{c.1.-checklist-tuxe1i-lux1eadp-trux1b0ux1edbc-khi-nux1ed9p}

\begin{itemize}
\item
  Lưu requirements.txt hoặc environment.yml kèm phiên bản Python và thư
  viện.
\item
  Lưu random seed, cách chia dữ liệu và hash của các file
  metadata/split.
\item
  Lưu model, StandardScaler, PCA và cấu hình feature extraction theo
  cùng một version.
\item
  Lưu command hoặc notebook dùng để tạo từng bảng kết quả trong báo cáo.
\item
  Chạy lại pipeline từ dữ liệu đầu vào đến metric cuối trên một môi
  trường sạch.
\item
  Đối chiếu số chiều HOG thực tế với artifact 1.812 trước khi công bố
  decomposition.
\end{itemize}

\end{document}
